% GENERATED FILE. Edit canonical lessons, then run npm run generate:books.
% canonical-source-hash: fnv1a64:f691104ef7218647
% canonical-lessons: ZH-C400-tiaozheng, ZH-C400-yanchang, ZH-C400-suoduan, ZH-C400-yuanding, ZH-C400-anding

\chapter{Moving the Dates}
\label{ch:zh-moving-the-dates}
\hlchaptermodality{400}

\begin{chapteropening}
\textbf{By the end of this chapter:} \emph{I can talk about adjusting a plan, making a stay longer or a trip shorter, the original plan and a settled life.}

Complete the last lesson of chapter 400: I can talk about adjusting a plan, making a stay longer or a trip shorter, the original plan and a settled life.
\end{chapteropening}

\section[tiáozhěng]{tiáozhěng — to adjust}

\label{lesson:ZH-C400-tiaozheng}



\begin{quote}
Before the new one: say the Mandarin for to meet (a requirement), then the Mandarin for to marry and settle down.
\end{quote}

\subsection*{You'll want to know: tiáozhěng}
\emph{tiáozhěng} — \textquotedblleft{}to adjust\textquotedblright{}.

\begin{quote}
\emph{rúguǒ shíjiān bù héshì wǒmen kěyǐ tiáozhěng} — If the time doesn't suit, we can adjust it.
\end{quote}

Its characters come later; for now it is a word for the ear.

\begin{cousinweb}
\emph{tiáozhěng} is built from \emph{tiáo} 'to put in tune' and \emph{zhěng} 'to put in order', both inherited from Middle Chinese, \emph{dew} and \emph{tsyengX} (Guangyun, in Baxter's transcription), which glosses the first 'to harmonise' and the second 'straight' and 'neat'.
\end{cousinweb}

\begin{grammarlens}[title={What you've built}]
One more word for saying what would happen.
\end{grammarlens}

\subsection*{Your turn}
\begin{itemize}
\raggedright
  \item \emph{Say it:} \emph{tiáozhěng}
  \item \emph{Say it:} \emph{tiáozhěng}, once more
  \item \emph{Say it:} \emph{rúguǒ shíjiān bù héshì wǒmen kěyǐ tiáozhěng} — If the time doesn't suit, we can adjust it.
  \item \emph{Recall:} say the Mandarin for to meet (a requirement), then the Mandarin for to marry and settle down, then say \emph{tiáozhěng} again
\end{itemize}

\begin{tcolorbox}[breakable,colback=teal!4,colframe=teal!35!black,title={Before you move on}]
What does tiáozhěng mean? (\textquotedblleft{}To adjust\textquotedblright{}.) Say it once more.
\end{tcolorbox}
\section[yáncháng]{yáncháng — to extend (in time)}

\label{lesson:ZH-C400-yanchang}



\begin{quote}
Before the new one: say the Mandarin for to marry and settle down, then the Mandarin for to adjust.
\end{quote}

\subsection*{You'll want to know: yáncháng}
\emph{yáncháng} — \textquotedblleft{}to extend (in time)\textquotedblright{}.

\begin{quote}
\emph{yàoshi nǐ xǐhuan zhèr wǒmen kěyǐ yáncháng jiàqī} — If you like it here, we can make the holiday longer.
\end{quote}

Its characters come later; for now it is a word for the ear.

\begin{cousinweb}
\emph{yáncháng} is built from \emph{yán} 'to stretch' and \zh{长} 'long', both inherited from Middle Chinese, \emph{yen} and \emph{drjang} (Guangyun, in Baxter's transcription), which glosses the first 'far' and 'to lengthen'.
\end{cousinweb}

\begin{grammarlens}[title={What you've built}]
One more word for saying what would happen.
\end{grammarlens}

\subsection*{Your turn}
\begin{itemize}
\raggedright
  \item \emph{Say it:} \emph{yáncháng}
  \item \emph{Say it:} \emph{yáncháng}, once more
  \item \emph{Say it:} \emph{yàoshi nǐ xǐhuan zhèr wǒmen kěyǐ yáncháng jiàqī} — If you like it here, we can make the holiday longer.
  \item \emph{Recall:} say the Mandarin for to marry and settle down, then the Mandarin for to adjust, then say \emph{yáncháng} again
\end{itemize}

\begin{tcolorbox}[breakable,colback=teal!4,colframe=teal!35!black,title={Before you move on}]
What does yáncháng mean? (\textquotedblleft{}To extend (in time)\textquotedblright{}.) Say it once more.
\end{tcolorbox}
\section[suōduǎn]{suōduǎn — to shorten}

\label{lesson:ZH-C400-suoduan}



\begin{quote}
Before the new one: say the Mandarin for to adjust, then the Mandarin for to extend (in time).
\end{quote}

\subsection*{You'll want to know: suōduǎn}
\emph{suōduǎn} — \textquotedblleft{}to shorten\textquotedblright{}.

\begin{quote}
\emph{zuò fēijī kěyǐ suōduǎn yíbàn de shíjiān} — Flying can cut the journey time by half.
\end{quote}

Its characters come later; for now it is a word for the ear.

\begin{cousinweb}
\emph{suōduǎn} is built from \emph{suō} 'to draw in' and \emph{duǎn} 'short', both inherited from Middle Chinese, \emph{srjuwk} and \emph{twanX} (Guangyun, in Baxter's transcription), which glosses the first 'to shrink' and the second 'not long'.
\end{cousinweb}

\begin{grammarlens}[title={What you've built}]
One more word for saying what would happen.
\end{grammarlens}

\subsection*{Your turn}
\begin{itemize}
\raggedright
  \item \emph{Say it:} \emph{suōduǎn}
  \item \emph{Say it:} \emph{suōduǎn}, once more
  \item \emph{Say it:} \emph{zuò fēijī kěyǐ suōduǎn yíbàn de shíjiān} — Flying can cut the journey time by half.
  \item \emph{Recall:} say the Mandarin for to adjust, then the Mandarin for to extend (in time), then say \emph{suōduǎn} again
\end{itemize}

\begin{tcolorbox}[breakable,colback=teal!4,colframe=teal!35!black,title={Before you move on}]
What does suōduǎn mean? (\textquotedblleft{}To shorten\textquotedblright{}.) Say it once more.
\end{tcolorbox}
\section[yuándìng]{yuándìng — originally planned}

\label{lesson:ZH-C400-yuanding}



\begin{quote}
Before the new one: say the Mandarin for to extend (in time), then the Mandarin for to shorten.
\end{quote}

\subsection*{You'll want to know: yuándìng}
\emph{yuándìng} — \textquotedblleft{}originally planned\textquotedblright{}.

\begin{quote}
\emph{yuándìng de shíjiān gǎi le} — The time we first planned has changed.
\end{quote}

Its characters come later; for now it is a word for the ear.

\begin{cousinweb}
\emph{yuándìng} is built from \emph{yuán} 'original' and \zh{定} 'settled', both inherited from Middle Chinese, \emph{ngjwon} and \emph{dengH} (Guangyun, in Baxter's transcription), which glosses the first 'a broad plain'. The sense 'original' is not given there.
\end{cousinweb}

\begin{grammarlens}[title={What you've built}]
One more word for saying what would happen.
\end{grammarlens}

\subsection*{Your turn}
\begin{itemize}
\raggedright
  \item \emph{Say it:} \emph{yuándìng}
  \item \emph{Say it:} \emph{yuándìng}, once more
  \item \emph{Say it:} \emph{yuándìng de shíjiān gǎi le} — The time we first planned has changed.
  \item \emph{Recall:} say the Mandarin for to extend (in time), then the Mandarin for to shorten, then say \emph{yuándìng} again
\end{itemize}

\begin{tcolorbox}[breakable,colback=teal!4,colframe=teal!35!black,title={Before you move on}]
What does yuándìng mean? (\textquotedblleft{}Originally planned\textquotedblright{}.) Say it once more.
\end{tcolorbox}
\section[\texorpdfstring{\zh{安定} (āndìng)}{安定 (āndìng)}]{\zh{安定} (āndìng) — settled; stable}

\label{lesson:ZH-C400-anding}



\begin{quote}
Before the new one: say the Mandarin for to shorten, then the Mandarin for originally planned.
\end{quote}

\subsection*{You'll want to know: \zh{安定}}
\textbf{\zh{安定}} — \emph{āndìng} — \textquotedblleft{}settled; stable\textquotedblright{}.

\begin{quote}
\textbf{\zh{只有工作安定了我们才能成家}} \emph{zhǐyǒu gōngzuò āndìng le wǒmen cái néng chéngjiā} — Only once our jobs are settled can we get married.
\end{quote}

\begin{cousinweb}
\textbf{\zh{安定}} (\emph{āndìng}) is built from \zh{安} 'calm' and \zh{定} 'settled', both inherited from Middle Chinese, \emph{'an} and \emph{dengH} (Guangyun, in Baxter's transcription), which glosses the second with the first: 'calm'.
\end{cousinweb}

\begin{grammarlens}[title={What you've built}]
One more word for saying what would happen.
\end{grammarlens}

\subsection*{Your turn}
\begin{itemize}
\raggedright
  \item \emph{Say it:} \emph{āndìng}
  \item \emph{Say it:} \emph{āndìng}, once more
  \item \emph{Say it:} \emph{zhǐyǒu gōngzuò āndìng le wǒmen cái néng chéngjiā} — Only once our jobs are settled can we get married.
  \item \emph{Recall:} say the Mandarin for to shorten, then the Mandarin for originally planned, then say \emph{āndìng} again
\end{itemize}

\begin{tcolorbox}[breakable,colback=teal!4,colframe=teal!35!black,title={Before you move on}]
What does \zh{安定} mean? (\textquotedblleft{}Settled; stable\textquotedblright{}.) Say it once more.
\end{tcolorbox}
